\documentclass[10pt,a4paper]{amsart}
\usepackage[margin=19mm]{geometry}
\usepackage{amsmath,amssymb,mathtools}
\usepackage[scr=rsfs,cal=euler]{mathalfa}
\usepackage{xfrac}
\usepackage[unicode,hidelinks,bookmarks=false]{hyperref}
\newcommand{\ie}{i.e.\ }
\newcommand{\cf}{cf.\ }
\newcommand{\resp}{respectively\ }
\newcommand{\Subsection}{\subsection}
\providecommand{\nobreakdash}{\nobreak\hbox{-}}
\setlength{\emergencystretch}{2em}
\multlinegap0pt
\usepackage{bm}
\usepackage{bbm}
\usepackage{dsfont}
\usepackage{xspace}
\usepackage{cancel}
\usepackage{mathtools}
\usepackage{amsthm}
\makeatletter
\@ifundefined{prop}{\newtheorem{prop}{Prop}}{}
\@ifundefined{lem}{\newtheorem{lem}[prop]{Lemma}}{}
\makeatother
\DeclareMathOperator {\Imm} {Im}
\newcommand {\p} {\partial}
\newcommand{\Om}{\Omega}
\newcommand{\dbar}{\ol{\p}}
\newcommand{\eps}{\epsilon}
\newcommand{\e}{\epsilon}
\newcommand{\norm}[1]{\lVert #1 \rVert}       
\newcommand{\ol}[1]{\overline{#1}}
\newcommand{\op}{\overline \p}
\newcommand{\oz}{\overline{z}}
\title[An Inverse Problem for the Prescribed Mean Curvature]{An Inverse Problem for the Prescribed Mean Curvature}
\author{Tony Liimatainen, Janne Nurminen}
\date{Source: arXiv:2509.22078v2 (CC BY 4.0)}
\begin{document}
\maketitle

\noindent\emph{Source and licence.} An Inverse Problem for the Prescribed Mean Curvature. Version arXiv:2509.22078v2 (CC BY 4.0), distributed under the Creative Commons Attribution 4.0 International licence, as recorded in the arXiv licence metadata for this exact version and shown on the abstract page https://arxiv.org/abs/2509.22078v2. Full rights notice: licenses/arxiv-2509.22078-CC-BY-4.0.txt.\par\smallskip

\noindent\textit{Editorial scope.} Counted unit: the Lem whose source label is \texttt{lemma\_D} in the article (source0.tex lines 1452--1454, source label \texttt{lemma\_D}), delivered together with the article's own complete original proof of it (source0.tex lines 1456--1528, one \texttt{proof} environment of 73 lines). No mathematics is written, repaired or re-derived: statement and proof are the article's own text line for line. Required context retained uncounted: eq:sobo\_decayL2 (lines 896--901); eq:rh\_form (lines 911--913); eq:sh\_formula (lines 920--922); eq:CZ\_rh\_norms (lines 935--937); eq: no critical point estimate (lines 955--958); eq: no critical point expansion II (lines 960--961); holomorp\_coord (lines 1233--1235). Every definition, display and label of those lines is retained unchanged. Omitted from this package and not redistributed: 1--895, 902--910, 914--919, 923--934, 938--954, 959--959, 962--1232, 1236--1451, 1455--1455, 1529--1715. Nothing inside a retained excerpt is omitted or altered: every retained body is delivered exactly as the article's lines are written, so no omission receipt of any kind accompanies this package. Citations: the retained excerpts carry no citation command at all, so no bibliography entry of the article is transcribed and no reference is redistributed. Reference adaptation: none was needed and none was made; every reference inside the delivered unit resolves inside it, so no unresolved reference or citation is produced. Printed slips kept literal: no printed word, symbol, hypothesis or number of the counted statement or of its proof was altered. Layout and class: the article's own class is \texttt{amsart} with packages including amsmath, amscd, amssymb, amsthm, comment, graphicx, epstopdf, mathrsfs, xcolor, hyperref; the wrapper is the lane's pinned \texttt{amsart} editorial wrapper at 10pt on A4, which restates the theorem-like environments the retained text needs and adds the article's own macro definitions that the retained text needs (\texttt{\textbackslash Imm}, \texttt{\textbackslash Om}, \texttt{\textbackslash dbar}, \texttt{\textbackslash e}, \texttt{\textbackslash eps}, \texttt{\textbackslash norm}, \texttt{\textbackslash ol}, \texttt{\textbackslash op}, \texttt{\textbackslash oz}, \texttt{\textbackslash p}), with extra wrapper packages (bm, bbm, dsfont, xspace, cancel, mathtools). The delivered numbering is the unit's own. Assets: no retained span contains a figure environment, a graphics-inclusion command, a PDF-page inclusion, a table or a TikZ picture; no image file is copied into this package, the package declares no asset dependency and no image file is redistributed with it. Licence: version arXiv:2509.22078v2 of this article is distributed under the Creative Commons Attribution 4.0 International licence, as recorded in the arXiv licence metadata for this exact version and shown on the abstract page https://arxiv.org/abs/2509.22078v2. A full rights notice is delivered at licenses/arxiv-2509.22078-CC-BY-4.0.txt. Characters: every retained excerpt is pure ASCII, so the delivered engine, XeLaTeX with its default Latin Modern text font, reports no missing character.
\begin{align}\label{eq:sobo_decayL2}
\begin{split}
 \norm{\overline \p_\psi^{-1}f}_{L^2(M)}&\leq C h^{1/2+\eps}\norm{f}_{W^{1,p}(M, T_{0,1}^*M)} %\\
 % \norm{\overline \p_\psi^{*-1}\omega}_{L^2(M)}&\leq C h^{1/2+\eps}\norm{\omega}_{W^{1,p}(M, T_{1,0}^*M)}.
 \end{split}
\end{align}

\begin{equation}\label{eq:rh_form}
 r_h=-  \overline{\p}_\psi^{-1}\sum_{j=0}^\infty T_h^j\op_\psi^{*-1}(qa),
\end{equation}

\begin{equation}\label{eq:sh_formula}
s_h=\sum_{j=0}^\infty T_h^j\op_\psi^{*-1}(qa)
\end{equation}

\begin{equation}\label{eq:CZ_rh_norms}
||r_h||_{L^p}, ||s_h||_{L^p},   ||\nabla r_h||_{L^p}=O(h^{1/p+\epsilon}),
\end{equation}

\begin{eqnarray}
\label{eq: no critical point estimate}
\|r_h\|_p +\|s_h\|_p+ \|d r_h\|_p \leq Ch
\end{eqnarray}

\begin{eqnarray}\label{eq: no critical point expansion II} \op_\psi^{-1} f =e^{-2i\psi/h} \frac{ih}{2} \frac{f}{\bar\partial\psi} + \frac{ih}{2}\op^{-1}\left( e^{-2i\psi/h}\bar\partial\left(\frac{f}{\bar\partial\psi}\right)\right),
\end{eqnarray} 

\begin{equation}\label{holomorp_coord}
	g(\nabla u,\nabla v)=2c^{-1}(\p u\dbar v + \dbar u\p v),\quad\Delta_gu = -2ic^{-1}\ol{\p}\p u\quad\text{and}\quad dV_g=c(z)\frac{i}{2}dzd\oz.
\end{equation}

\begin{lem}\label{lemma_D}
	For $I_1$ we have $\lim_{h\to 0} hI_1=0.$
\end{lem}

\begin{proof}
	We immediately use \eqref{holomorp_coord} to have
	\begin{align*}
		&\sum_{k=1}^{2} \int_{\Om} T^k\p_kv_0g(\nabla v_1, \nabla v_2)\,dV_g\\
		&= \sum_{k=1}^2 i \int_{\Om} T^k\left(i^{k-1}\p v_0 + (-i)^{k-1}\dbar v_0\right)\left(\p v_1 \dbar v_2 + \dbar v_1\p v_2\right)\,dzd\oz\\
		&= \sum_{k=1}^2 i \int_{\Om} D_1\p v_0 \p v_1\dbar v_2 + D_2\p v_0 \dbar v_1\p v_2 + D_3\dbar v_0 \p v_1\dbar v_2 + D_4\dbar v_0 \dbar v_1\p v_2\,dzd\oz.
	\end{align*}
	Here all $D_j$ again depend on the summation index $k$.
	
	Starting with the $D_1$ term we expand further (and as before we do not write the summation):
	\begin{align}\label{D_1_tilde}
		&\int_{\Om} D_1\p v_0 \p v_1\dbar v_2\,dzd\oz\\\notag
		&= \int_{\Om}\tilde D_1 \hat{v}_0\hat{v}_1\hat{v}_2 + \tilde D_1\hat{v}_0\hat{v}_1\dbar \hat{v}_2 + \tilde D_1\hat{v}_0\p \hat{v}_1 \hat{v}_2 + \tilde D_1\hat{v}_0\p \hat{v}_1\dbar \hat{v}_2 + \tilde D_1 \p \hat{v}_0\hat{v}_1\dbar \hat{v}_2 \\\notag
		&+ \tilde D_1\p \hat{v}_0\p \hat{v}_1 \hat{v}_2 + \tilde D_1 \p \hat{v}_0\hat{v}_1\hat{v}_2 + \tilde D_1\p \hat{v}_0 \p \hat{v}_1\dbar \hat{v}_2\,dzd\oz.
	\end{align}
	Here we have abused notation using $\tilde D_1$ for all the terms that do not depend on $h$. The first six terms are similar to $A_1, A_3, A_2, B_5, B_4$ and $B_3$, in that order, and they are $O(h^{-1/2+\e})$. For the first remaining term we have
	\begin{align*}
		&\int_{\Om}\tilde D_1 \p \hat{v}_0\hat{v}_1\hat{v}_2\,dzd\oz\\
		&= \int_{\Om}\tilde D_1 \p \hat{v}_0 \hat{v}_1\hat{v}_2\,dzd\oz\\
		&=-\frac{2}{h}\int_{\Om}\tilde D_1 \p\ol{\Phi} e^{-2\ol{\Phi}/h}(a_0+r_0)\hat{v}_1\hat{v}_2 + \tilde D_1e^{(2\Phi - 2\ol{\Phi})/h}\p r_0(a_1 + r_1)(a_2 + r_2)\,dzd\oz.
	\end{align*}
	The first term is $O(1)$ by stationary phase and the fact that $\p\ol{\Phi}(z_0)=0$ and the second term is $O(h^{1/2+\e})$ because of the term $\p r_0$ (see \eqref{eq:CZ_rh_norms}). For the last term in \eqref{D_1_tilde} we use that $\p a_0 =\dbar a_2 = \dbar(-\Psi + \Phi)=0=\dbar\ol{\Phi}(z_0)$ and stationary phase to have
	\begin{align*}
		&\int_{\Om}\tilde D_1\p \hat{v}_0\p \hat{v}_1\dbar \hat{v}_2\,dzd\oz\\
		&= \int_{\Om}e^{(2\Phi-2\ol{\Phi})/h}\p r_0\left(\frac{1}{h}\p(\Psi+\Phi)(a_1+r_1) + \p a_1 + \p r_1\right)\p r_2\,dzd\oz\\
		&= O(h^{1/2+\e}).
	\end{align*}
	In the last equality we used the estimates for $\p r_0, \p r_2$ (see \eqref{eq:CZ_rh_norms}, \eqref{eq: no critical point estimate}).
	
	The term with $D_2$ goes similarly as above with the term $D_1$ because $v_1$ and $v_2$ have holomorphic phases.
	
	Moving to the term with $D_3$ we have (using similar notation as with $D_1$)
	\begin{align*}
		&\int_{\Om}D_3\dbar v_0 \p v_1\dbar v_2\,dzd\oz\\
		&= \int_{\Om}\tilde D_3\hat{v}_0\hat{v}_1\hat{v}_2 + \tilde D_3\hat{v}_0\hat{v}_1\dbar \hat{v}_2 + \tilde D_3\hat{v}_0\p \hat{v}_1\hat{v}_2 + \tilde D_3\hat{v}_0\p \hat{v}_1\dbar \hat{v}_2 + \tilde D_3\dbar \hat{v}_0\hat{v}_1\dbar \hat{v}_2\\
		&+ \tilde D_3\dbar \hat{v}_0\hat{v}_1\hat{v}_2 + \tilde D_3\dbar \hat{v}_0\p \hat{v}_1\dbar \hat{v}_2\,dzd\oz.
	\end{align*}
	Again the first six terms are similar to ones before. These are $O(h^{-1/2+\e})$ by the arguments made for $A_1, A_3, A_2, B_5, B_6$ and $B_1$. The first remaining term is
	\begin{align*}
		&\int_{\Om}\tilde D_3\dbar \hat{v}_0 \hat{v}_1\hat{v}_2\,dzd\oz\\
		&= \int_{\Om}e^{(2\Phi-2\ol{\Phi})/h}\left(\left(-\frac{2}{h}\dbar \ol{\Phi}(a_0 + r_0) + \dbar a_0 + \dbar r_0\right)(a_1+r_1)(a_2+r_2)\right)\,dzd\oz\\
		&= O(1)
	\end{align*}
	by stationary phase and the estimates for $r_k$ (see \eqref{eq:CZ_rh_norms}, \eqref{eq: no critical point estimate}). The last remaining term is, by using again that the functions $a_2, \Phi$ and $\Psi$ are holomorphic,
	\begin{align*}
		&\int_{\Om}\tilde D_3 \dbar \hat{v}_0\p \hat{v}_1\dbar \hat{v}_2\,dzd\oz\\
		&= \int_{\Om}\tilde D_3 e^{(2\Phi -2\ol{\Phi})/h}\left(-\frac{2}{h}\dbar\ol{\Phi}(a_0+r_0) + \dbar a_0 + \dbar r_0\right)\\
		&\times\left(\frac{1}{h}\p(\Psi+\Phi)(a_1+r_1) + \p a_1 + \p r_1\right)\dbar r_2\,dzd\oz\\
		&= -\frac{2}{h^2}\int_{\Om}\tilde D_3 e^{(2\Phi -2\ol{\Phi})/h}\dbar\ol{\Phi}\p(\Psi+\Phi)a_1a_0\dbar r_2\,dzd\oz + O(h^{-1/2+\e})
	\end{align*}
	because of the estimates for $r_1, r_0$ and $\dbar r_2$ (see \eqref{eq:CZ_rh_norms}, \eqref{eq: no critical point estimate}). For the last term, we recall that (see \eqref{eq:rh_form} - \eqref{eq:sh_formula})
	\begin{align*}
		r_2 &= - \dbar^{-1}_{\Imm(-\Psi + \Phi)}s_2 = \dbar^{-1}e^{i\Imm(-\Psi + \Phi)/h}s_2\\
		s_2 &= -\op_{\Imm(-\Psi + \Phi)}^{*-1}(qa) - \sum_{j=1}^{\infty}T_2^j\op_{\Imm(-\Psi + \Phi)}^{*-1}(qa)
	\end{align*}
	and thus in holomorphic coordinates, denoting $\eta= \Imm(-\Psi + \Phi)$,
	\begin{align*}
		\dbar r_2 &= -e^{i\eta/h}s_2\\
		&= e^{i\eta/h}\p^{-1}(e^{i\eta/h}qa) + e^{i\eta/h}\sum_{j=1}^{\infty}T_2^j\p^{-1}(e^{i\eta/h}qa)\\
		&= e^{i\eta/h}\left(\frac{ih}{2}e^{i\eta/h}\frac{qa}{\p\eta} + \frac{ih}{2}\p^{-1}\left(e^{i\eta/h}\p\left(\frac{qa}{\p\eta}\right)\right)\right) + e^{i\eta/h}\sum_{j=1}^{\infty}T_2^j\p^{-1}(e^{\eta/h}qa).
	\end{align*}
    Here we used the expansion \eqref{eq: no critical point expansion II}.
	Now
	\begin{align*}
		-\frac{2}{h^2}\int_{\Om}\tilde D_3 e^{(2\Phi -2\ol{\Phi})/h}\dbar\ol{\Phi}\p(\Psi+\Phi)a_1a_0\frac{ih}{2}e^{2i\eta/h}\frac{qa}{\p\eta}\,dzd\oz = O(1)
	\end{align*}
	by stationary phase.
	Recall that 
    $$\p^{-1}\left(e^{i\eta/h}\p\left(\frac{qa}{\p\eta}\right)\right) = O(h^{1/2+\e}),\quad \p^{-1}(e^{i\eta/h}qa) = O(h)$$
    (by the estimate \eqref{eq:sobo_decayL2} and the expansion \eqref{eq: no critical point expansion II}) and $T_h = O(h^{1/2-\e})$. Then the remaining terms in $\int_{\Om}\tilde D_3 \dbar \hat{v}_0\p \hat{v}_1\dbar \hat{v}_2\,dzd\oz$ are $O(h^{3/2-\e})$
	
	Again the term with $D_4$ goes similarly as above with the term $D_3$ because $v_1$ and $v_2$ have holomorphic phases.
\end{proof}

\end{document}
